\documentclass[11pt]{article}
\usepackage{amsmath,amssymb}
\usepackage{graphicx}
\usepackage{hyperref}
\title{LSC 6.2: A Phenomenological Framework for Neutrino Propagation and Detector-Frame Tensor Anisotropy}
\author{LuciferSun \\ Independent Researcher}
\date{}
\begin{document}
\maketitle
\begin{abstract}
LSC 6.2 is an arXiv-safe synthesis draft extending the public LSC 6.0 phenomenological framework. It keeps the conservative propagation-measurement coupling of LSC 6.0 and formulates detector-frame anisotropy through an explicit rank-2 tensor. The model is not presented as confirmed new physics; it is an effective proposal requiring global fitting, numerical simulation, and independent constraints from gallium experiments, KATRIN, and IceCube.
This work does not claim a fundamental origin of the effect, but provides a testable phenomenological ansatz.
\end{abstract}

\section{Introduction}
LSC 6.2 is a conservative extension draft of the public LSC 6.0 framework. The public LSC 6.0 record is available at \url{https://zenodo.org/records/19780616}, and the related GitHub release is available at \url{https://github.com/luciferprosun/LSC-6.0/releases/tag/v1.0.0}. We propose a detector-frame tensor ansatz for neutrino energy reconstruction and angular response.

The effect can be interpreted as an effective direction- and energy-dependent phase correction, without invoking a literal refractive index for neutrinos. The model remains phenomenological.

\section{Theoretical Framework}
The core LSC 6.0 observable ansatz is retained,
\begin{equation}
 E_{\rm true}(x)=E_\infty\,G(g_{\mu\nu},\Phi(x)),
\end{equation}
\begin{equation}
 E_{\rm rec}=E_{\rm true}\left(1+\alpha_D D_{ij}\hat p^i\hat p^j\right).
\end{equation}
Here $G$ is a dimensionless propagation factor, $E_\infty$ is the conserved reference energy, $D_{ij}$ is an effective spatial detector-response tensor in the laboratory frame, $\hat p^i$ is the reconstructed direction unit vector, and $\alpha_D$ controls the detector-level contribution. The limiting condition is
\begin{equation}
 \lim_{\Phi\to 0,\,\alpha_D\to 0}E_{\rm rec}=E_\infty.
\end{equation}

\section{Mathematical Formulation}
The propagation factor is expanded as
\begin{equation}
 G(g_{\mu\nu},\Phi(x))=1+\delta_G(x)+O(\delta_G^2).
\end{equation}
The detector tensor is decomposed into isotropic and anisotropic components,
\begin{equation}
 D_{ij}=D^{\rm iso}_{ij}+\Delta D_{ij}.
\end{equation}
The effective fractional shift is then
\begin{equation}
 \frac{\Delta E}{E}\simeq \delta_G+\alpha_D\Delta_D,
 \qquad
 \Delta_D=D_{ij}\hat p^i\hat p^j.
\end{equation}

The anisotropic part is parameterized by a traceless rank-2 tensor,
\begin{equation}
 D_{ij}=\delta\left(n_i n_j-\frac{1}{3}\delta_{ij}\right),
\end{equation}
where $\delta$ is the dimensionless anisotropy magnitude and $n_i$ is a preferred detector-frame direction. For an IceCube-style implementation, $n_i$ may be aligned with a measured ice-flow or crystal-orientation-fabric axis; the corresponding ice-light propagation anisotropy is treated as a detector systematic, not automatically as a neutrino-sector signal.
The contraction entering the observable response is
\begin{equation}
 D_{ij}\hat p^i\hat p^j=
 \delta\left[(\hat p\cdot n)^2-\frac{1}{3}\right].
\end{equation}

The reconstructed energy is therefore modeled as
\begin{equation}
 E_{\rm rec}(\hat p)=E_{\rm true}\left(1+\alpha_D D_{ij}\hat p^i\hat p^j\right).
\end{equation}
A bounded first-order effect may be summarized as
\begin{equation}
 \Delta m_{ij,{\rm eff}}^2\simeq
 \Delta m_{ij}^2\left(1+\epsilon_D\right),
 \qquad
 \epsilon_D=D_{kl}\hat p^k\hat p^l .
\end{equation}

The effective energy entering the oscillation phase is
\begin{equation}
 E_{\rm eff}(x)=E_{\rm true}(x)
 \left[1+\alpha_D D_{ij}\hat p^i\hat p^j\right].
\end{equation}
The oscillation phase is then written in the standard ultra-relativistic form
\begin{equation}
 \Delta \Phi_{ij}=
 \int_0^L \frac{\Delta m_{ij}^2}{2E_{\rm eff}(x)}\,dx .
\end{equation}
This is a compact phenomenological ansatz and not a microscopic derivation.
In the limit $D_{ij}\to 0$ and $\delta_G\to 0$, the standard three-flavor neutrino oscillation framework is recovered.

\section{Phenomenological Implications}
The relevant output of LSC 6.2 is a set of test classes: angular anisotropy, detector dependence, energy-dependent residuals, and sidereal modulation. Representative archive diagnostics are shown in Figs.~\ref{fig:lsc62-framework}--\ref{fig:lsc62-katrin}.
A representative value $\delta\sim 0.05$ is used for illustration; it is not a fitted or claimed measurement.

For binned event data, the conservative comparison can be written as
\begin{equation}
 \chi^2=\sum_k
 \frac{\left[N_{{\rm obs},k}-N_{{\rm pred},k}(D_{ij},\delta_G)\right]^2}
 {\sigma_k^2},
\end{equation}
with
\begin{equation}
 \sigma_k^2=\sigma_{{\rm stat},k}^2+\sigma_{{\rm sys},k}^2.
\end{equation}
Each bin $k$ may include reconstructed energy, zenith angle, azimuth, and time. A null test corresponds to $D_{ij}\rightarrow 0$ and $\delta_G\rightarrow 0$ within experimental uncertainties.

\begin{figure}[ht]
\centering
\includegraphics[width=0.82\linewidth]{figures/fig1_framework.png}
\caption{Unified propagation and detector-response framework.}
\label{fig:lsc62-framework}
\end{figure}

\begin{figure}[ht]
\centering
\includegraphics[width=0.82\linewidth]{figures/fig2_anisotropy.png}
\caption{Anisotropic response diagnostic for angular or directional effects.}
\label{fig:lsc62-anisotropy}
\end{figure}

\begin{figure}[ht]
\centering
\includegraphics[width=0.82\linewidth]{figures/fig3_parameter_scan.png}
\caption{Parameter-scan diagnostic for effective LSC coefficients.}
\label{fig:lsc62-scan}
\end{figure}

\begin{figure}[ht]
\centering
\includegraphics[width=0.82\linewidth]{figures/fig4_katrin.png}
\caption{KATRIN-style consistency diagnostic.}
\label{fig:lsc62-katrin}
\end{figure}

\section{Discussion}
The tensor language is the primary mathematical structure of LSC 6.2. The effect can be interpreted as an effective direction- and energy-dependent phase correction, without invoking a literal refractive index for neutrinos. The model must be checked against gallium source experiments, KATRIN beta-spectrum constraints, IceCube anisotropy searches, and standard three-flavor oscillation data. Known IceCube ice-light propagation anisotropies, including crystal-orientation fabric and hole-ice effects, must be treated as backgrounds or nuisance systematics before any neutrino-sector interpretation is attempted.

\section{Conclusion}
LSC 6.2 provides a stable arXiv-safe synthesis of the LSC 6.0 effective model with an explicit detector-frame anisotropy tensor. The next required step is numerical fitting with explicit priors on $\delta_G$, $\alpha_D$, and the angular coefficients of $D_{ij}$.

\section{References}
\begin{thebibliography}{9}
\bibitem{wolfenstein1978} L. Wolfenstein, Neutrino oscillations in matter, Phys. Rev. D 17, 2369 (1978).
\bibitem{msw1985} S. P. Mikheyev and A. Yu. Smirnov, Resonant amplification of neutrino oscillations in matter and solar neutrino spectroscopy, Sov. J. Nucl. Phys. 42, 913 (1985).
\bibitem{best2022} V. V. Barinov et al., Results from the Baksan Experiment on Sterile Transitions (BEST), Phys. Rev. Lett. 128, 232501 (2022).
\bibitem{katrin} KATRIN Collaboration, Direct neutrino-mass constraints from tritium beta decay.
\bibitem{icecube} IceCube Collaboration, searches for Lorentz violation and anisotropy in high-energy neutrinos.
\bibitem{pbh} B. J. Carr and S. W. Hawking, Black holes in the early Universe, Mon. Not. R. Astron. Soc. 168, 399 (1974).
\bibitem{zenodo42} LuciferSun, LSC 4.2 ULTRA public record, Zenodo, \url{https://doi.org/10.5281/zenodo.19602045}.
\bibitem{zenodo60} LuciferSun, LSC 6.0 public record, Zenodo, \url{https://zenodo.org/records/19780616}.
\bibitem{github60} LuciferSun, LSC 6.0 GitHub release, \url{https://github.com/luciferprosun/LSC-6.0/releases/tag/v1.0.0}.
\end{thebibliography}

\end{document}
